\documentclass[11pt,a4paper]{article}
\usepackage{amsmath,amssymb,bm}
\usepackage[margin=2.5cm]{geometry}
\usepackage{booktabs}
\usepackage{hyperref}

\title{Electrohydrodynamic Simulation Theory\\
       \large VSEPR-SIM --- Coupled Navier-Stokes / Poisson / Nernst-Planck\\
       for Structured Channel Geometries}
\author{VSEPR-SIM Project}
\date{\today}

\begin{document}
\maketitle

% ============================================================================
\section{Scope}
% ============================================================================

This document establishes the theoretical foundation for the
\texttt{sim/ehd} module of VSEPR-SIM.  The module solves a coupled
multiphysics problem: incompressible fluid flow, electrostatics, and ionic
species transport inside structured (helical-electrode) channel geometries.

The governing equations are:
\begin{enumerate}
  \item \textbf{Continuity and Momentum} (incompressible Navier--Stokes)
  \item \textbf{Poisson equation} for the electrostatic potential
  \item \textbf{Nernst--Planck equation} for ionic species transport
\end{enumerate}

% ============================================================================
\section{Governing Equations}
% ============================================================================

\subsection{Fluid Flow}

Incompressible continuity:
\begin{equation}
  \nabla\cdot\bm{u} = 0
\end{equation}

Momentum conservation with electric body force:
\begin{equation}
  \rho\!\left(\frac{\partial\bm{u}}{\partial t}
  + \bm{u}\cdot\nabla\bm{u}\right)
  = -\nabla p + \mu\,\nabla^{2}\bm{u} + \bm{f}_{\mathrm{elec}}
\end{equation}

where the electric body force is
\begin{equation}
  \bm{f}_{\mathrm{elec}} = \rho_{e}\,\bm{E}
  = \rho_{e}\!\left(-\nabla\phi\right)
\end{equation}

and $\rho_{e}$ is the local space-charge density from ionic species.

Reynolds number for a tube of hydraulic diameter~$D_h$:
\begin{equation}
  \mathrm{Re} = \frac{\rho\,U\,D_h}{\mu}
\end{equation}

\subsection{Electrostatics}

Poisson equation for the electric potential~$\phi$:
\begin{equation}
  \nabla\cdot\!\left(\epsilon\,\nabla\phi\right) = -\rho_{e}
\end{equation}

Electric field:
\begin{equation}
  \bm{E} = -\nabla\phi
\end{equation}

Boundary conditions:
\begin{align}
  \phi &= V_{+} \quad\text{on electrode}_{+} \\
  \phi &= V_{-} \quad\text{on electrode}_{-} \\
  \bm{n}\cdot(\epsilon\,\nabla\phi) &= 0
    \quad\text{on insulating walls}
\end{align}

\subsection{Ion Transport (Nernst--Planck)}

Species conservation for ion~$i$:
\begin{equation}
  \frac{\partial c_{i}}{\partial t} + \nabla\cdot\bm{N}_{i} = 0
\end{equation}

Molar flux:
\begin{equation}\label{eq:NP}
  \bm{N}_{i}
  = \underbrace{-D_{i}\,\nabla c_{i}}_{\text{diffusion}}
    \underbrace{-\,z_{i}\,\mu_{i}\,F\,c_{i}\,\nabla\phi}_{\text{electromigration}}
    + \underbrace{c_{i}\,\bm{u}}_{\text{convection}}
\end{equation}

where $D_i$ is diffusivity, $z_i$ is valence, $\mu_i$ is ionic mobility, and
$F$ is Faraday's constant.

Nernst--Einstein relation (when mobility is not independently known):
\begin{equation}
  \mu_{i} = \frac{|z_i|\,F\,D_i}{R\,T}
\end{equation}

Space-charge density assembly:
\begin{equation}
  \rho_{e} = F\sum_{i} z_{i}\,c_{i}
\end{equation}

% ============================================================================
\section{Coupling Strategy}
% ============================================================================

A segregated outer-iteration scheme is used:
\begin{enumerate}
  \item Solve flow: $\bm{u}, p$ (Navier--Stokes with current $\bm{f}_{\mathrm{elec}}$)
  \item Solve electric: $\phi, \bm{E}$ (Poisson with current $\rho_e$)
  \item Solve species: $c_i, \bm{N}_i$ (Nernst--Planck with current $\bm{u}, \bm{E}$)
  \item Update $\rho_e = F\sum z_i c_i$
  \item Repeat until convergence
\end{enumerate}

Under-relaxation factors $\alpha_{\mathrm{flow}}, \alpha_{\mathrm{elec}},
\alpha_{\mathrm{species}}$ are applied to stabilise the outer loop.

% ============================================================================
\section{Analytical Baselines}
% ============================================================================

\subsection{Hagen--Poiseuille Flow (Straight Tube)}

Axial velocity profile:
\begin{equation}
  u_z(r) = 2\,\bar{U}\!\left(1 - \frac{r^2}{R^2}\right)
\end{equation}

Pressure drop:
\begin{equation}
  \Delta P = \frac{8\,\mu\,\bar{U}\,L}{R^2}
\end{equation}

\subsection{Coaxial Capacitor Field}

Potential between inner radius~$R_{\mathrm{in}}$ and outer radius~$R_{\mathrm{out}}$:
\begin{equation}
  \phi(r) = V\,\frac{\ln(R_{\mathrm{out}}/r)}{\ln(R_{\mathrm{out}}/R_{\mathrm{in}})}
\end{equation}

Radial field:
\begin{equation}
  E_r(r) = \frac{V}{r\,\ln(R_{\mathrm{out}}/R_{\mathrm{in}})}
\end{equation}

% ============================================================================
\section{Parametric Description}
% ============================================================================

The full parameter set for a helical EHD reactor case is:
\begin{equation}
  \mathcal{P} = \bigl\{
    R_{\mathrm{tube}},\;L,\;p_{\mathrm{helix}},\;d_{\mathrm{wire}},\;
    n_{\mathrm{turns}},\;g,\;V,\;Q,\;\mu,\;\epsilon,\;D_i,\;z_i
  \bigr\}
\end{equation}

Derived quantities:
\begin{align}
  R_{\mathrm{helix}} &= R_{\mathrm{tube}} - g - d_{\mathrm{wire}}/2 \\
  D_h &= 2\,R_{\mathrm{tube}} \\
  \mathrm{Re} &= \rho\,U\,D_h / \mu
\end{align}

% ============================================================================
\section{Engineering Metrics}
% ============================================================================

\begin{align}
  \Delta P &= P_{\mathrm{in}} - P_{\mathrm{out}} \\
  \bar{E}  &= \frac{1}{V_{\Omega}}\int_{\Omega}|\bm{E}|\,dV \\
  \dot{N}_i &= \int_{A_{\mathrm{out}}}\bm{N}_i\cdot\bm{n}\,dA
\end{align}

Enhancement factor between geometries:
\begin{equation}
  \eta = \frac{\text{performance (modulated geometry)}}
              {\text{performance (baseline)}}
\end{equation}

Weighted optimisation objective:
\begin{equation}
  J = w_1\,\bar{E} + w_2\,\dot{N}_{\mathrm{out}} - w_3\,\Delta P
\end{equation}

% ============================================================================
\section{Implementation Mapping}
% ============================================================================

\begin{table}[h]
\centering
\caption{EHD module stage structure}
\begin{tabular}{@{}lll@{}}
\toprule
Stage & Namespace & Key Classes \\
\midrule
1. CAD       & \texttt{vsepr::ehd::cad}     & HelixGenerator, TubeBodyGenerator, ElectrodeLayout, StepManifest \\
2. Domain    & \texttt{vsepr::ehd::domain}   & DomainExtract, NamedRegionRegistry, BCTemplate \\
3. Physics   & \texttt{vsepr::ehd::physics}  & FlowModel, ElectrostaticModel, IonTransportModel, CoupledSolver \\
4. Mesh      & \texttt{vsepr::ehd::mesh}     & MeshControlSpec, InflationSpec, RefinementZone \\
5. Post      & \texttt{vsepr::ehd::post}     & ContourExport, StreamlineExport, SectionProbe, ComparisonTable \\
\bottomrule
\end{tabular}
\end{table}

% ############################################################################
%   EXPANDED THEORY — Phases 1–6
%   Appended sections providing full mathematical treatment of each stage
%   and the cross-cutting dimensional / sensitivity framework.
% ############################################################################

% ============================================================================
\section{Phase 1 — Parametric Geometry and Helical Electrode Theory}
\label{sec:phase1}
% ============================================================================

The CAD stage generates the physical geometry of the helical-electrode EHD
reactor from a minimal parameter set~$\mathcal{P}$.  Every geometric body is
derived deterministically; no stochastic elements are introduced.

\subsection{Helix Centreline}

The electrode wire follows a circular helix of radius~$R_h$ and pitch~$p$
wound along the tube axis~$\hat{\bm{z}}$.  The centreline is parameterised by
the azimuthal angle~$\theta$:
\begin{equation}\label{eq:helix_centre}
  \bm{r}_h(\theta) = \bigl(R_h\cos\theta,\;
                            R_h\sin\theta,\;
                            z_0 + p\,\theta/(2\pi)\bigr)
\end{equation}
where $z_0$ is the axial offset of the first turn (the inlet development
length) and the total number of discrete points for $n_t$ turns at angular
resolution~$\Delta\theta$ is
\begin{equation}\label{eq:helix_npts}
  N_{\mathrm{pts}} = n_t \cdot (2\pi/\Delta\theta) + 1
\end{equation}

The helix radius is derived from the run-card geometry:
\begin{equation}\label{eq:helix_radius}
  R_h = R_{\mathrm{tube}} - g - d_w / 2
\end{equation}
where $g$ is the wall-clearance gap and $d_w$ is the wire diameter.

\subsection{Wire Cross-Section}

At each centreline station the wire occupies a disk of radius~$r_w = d_w/2$
in a plane normal to the local tangent.  The tangent vector is
\begin{equation}\label{eq:helix_tangent}
  \bm{t}(\theta) = \frac{d\bm{r}_h}{d\theta}\bigg/
                    \left|\frac{d\bm{r}_h}{d\theta}\right|
  ,\qquad
  \frac{d\bm{r}_h}{d\theta}
  = \bigl(-R_h\sin\theta,\; R_h\cos\theta,\; p/(2\pi)\bigr)
\end{equation}

A local orthonormal frame $(\bm{t}, \bm{n}, \bm{b})$ is constructed via a
reference-vector projection.  The cross-section circle is then
\begin{equation}\label{eq:wire_circle}
  \bm{r}_{\mathrm{wire}}(\alpha) = \bm{r}_h
    + r_w\bigl(\bm{n}\cos\alpha + \bm{b}\sin\alpha\bigr),
  \quad \alpha \in [0, 2\pi)
\end{equation}

\subsection{Tube Profile — Straight and Modulated}

The inner wall of the tube is a surface of revolution.  For the straight
geometry the radius is constant:
\begin{equation}\label{eq:tube_straight}
  R_{\mathrm{inner}}(z) = R_{\mathrm{tube}}
\end{equation}

For the modulated (sinusoidal bulge) geometry:
\begin{equation}\label{eq:tube_mod}
  R_{\mathrm{inner}}(z)
  = R_{\mathrm{tube}}
    + A_m\,\sin\!\left(\frac{2\pi\,z}{\lambda_m}\right)
\end{equation}
where $A_m$ is the modulation half-amplitude and $\lambda_m$ is the spatial
wavelength.  The outer wall is offset by the wall thickness~$t_w$:
\begin{equation}\label{eq:tube_outer}
  R_{\mathrm{outer}}(z) = R_{\mathrm{inner}}(z) + t_w
\end{equation}

Annular ring tessellation at axial station~$z$ produces a vertex set
\begin{equation}\label{eq:ring}
  \bm{r}_{\mathrm{ring}}(\theta_j, z) =
    \bigl(R\cos\theta_j,\; R\sin\theta_j,\; z\bigr),
  \qquad \theta_j = \frac{2\pi\,j}{N_\theta}
\end{equation}
for $j = 0, \ldots, N_\theta - 1$.

\subsection{Electrode Layout and STEP Manifest}

Each electrode is described by a polarity~$\sigma\in\{+, -, \mathrm{gnd}\}$,
an applied voltage, and a geometry type (helical sweep, coaxial cylinder, or
planar pair).  The layout assembles all electrodes into a device description
from which a STEP-style manifest of named bodies is generated.

The total wire length for a helical electrode is
\begin{equation}\label{eq:wire_length}
  L_{\mathrm{wire}} = n_t\,\sqrt{(2\pi R_h)^2 + p^2}
\end{equation}
and the approximate wire volume (used for domain subtraction) is
\begin{equation}\label{eq:wire_volume}
  V_{\mathrm{wire}} = \pi\,r_w^2\,L_{\mathrm{wire}}
\end{equation}

% ============================================================================
\section{Phase 2 — Domain Decomposition and Boundary Classification}
\label{sec:phase2}
% ============================================================================

The computational domain is partitioned into disjoint material regions:
\begin{equation}\label{eq:domain_decomp}
  \Omega = \Omega_f \;\cup\; \Omega_{s,c} \;\cup\; \Omega_{s,d}
\end{equation}
where $\Omega_f$ is the fluid region, $\Omega_{s,c}$ the solid conductors
(electrodes), and $\Omega_{s,d}$ the solid dielectrics (tube wall, support
structures).

\subsection{Fluid Volume Extraction}

The net fluid volume is the cylindrical envelope minus all embedded solid
volumes:
\begin{equation}\label{eq:fluid_volume}
  |\Omega_f| = \pi\,R_{\mathrm{tube}}^2\,L_{\mathrm{total}}
             - \sum_k V_{s,k}
\end{equation}
where $L_{\mathrm{total}} = L_{\mathrm{inlet}} + L + L_{\mathrm{outlet}}$ and
$V_{s,k}$ is the volume of the $k$-th embedded solid.

\subsection{Boundary Surface Taxonomy}

Every surface of $\partial\Omega_f$ is assigned a unique tag from a fixed
taxonomy:
\begin{align}
  \Gamma_{\mathrm{in}}     &\quad\text{— inlet face}  \\
  \Gamma_{\mathrm{out}}    &\quad\text{— outlet face}  \\
  \Gamma_{\mathrm{wall}}   &\quad\text{— no-slip insulating wall}  \\
  \Gamma_{\mathrm{elec}+}  &\quad\text{— positive electrode surface}  \\
  \Gamma_{\mathrm{elec}-}  &\quad\text{— negative / grounded electrode surface}
\end{align}
Each tag maps to a triple of physics-specific boundary conditions
$(\text{flow BC},\;\text{electric BC},\;\text{transport BC})$ as follows:

\begin{table}[h]
\centering
\caption{Boundary condition mapping per surface tag}
\begin{tabular}{@{}lccc@{}}
\toprule
Surface            & Flow         & Electric            & Transport \\
\midrule
$\Gamma_{\mathrm{in}}$    & Velocity inlet $\bar{U}$  & Neumann $\bm{n}\!\cdot\!\nabla\phi=0$ & Fixed $c_{i,0}$ \\
$\Gamma_{\mathrm{out}}$   & Pressure outlet $P_0$     & Neumann insulating  & Convective outflow \\
$\Gamma_{\mathrm{wall}}$  & No-slip $\bm{u}=\bm{0}$  & Neumann insulating  & Zero flux $\bm{N}_i\!\cdot\!\bm{n}=0$ \\
$\Gamma_{\mathrm{elec}+}$ & No-slip                   & Dirichlet $\phi=V_+$ & Zero flux \\
$\Gamma_{\mathrm{elec}-}$ & No-slip                   & Dirichlet $\phi=V_-$ & Zero flux \\
\bottomrule
\end{tabular}
\end{table}

\subsection{Multi-Species Boundary Conditions}

For $N_s$ ionic species, the inlet boundary assigns a fixed concentration
for each species:
\begin{equation}\label{eq:bc_species}
  c_i\big|_{\Gamma_{\mathrm{in}}} = c_{i,0}, \qquad i = 1,\ldots,N_s
\end{equation}
All insulating walls and electrode surfaces enforce zero normal flux:
\begin{equation}\label{eq:bc_zero_flux}
  \bm{N}_i \cdot \bm{n}\big|_{\Gamma_{\mathrm{wall}}\cup\Gamma_{\mathrm{elec}}} = 0
\end{equation}

% ============================================================================
\section{Phase 3 — Coupled Multiphysics: Discretisation and Solution Strategy}
\label{sec:phase3}
% ============================================================================

This section extends the governing equations from \S2 with the discrete
field representations, the finite-difference operators, and the segregated
coupling algorithm implemented in the \texttt{CoupledSolver}.

\subsection{Axisymmetric Structured Grid}

All three field variables are stored on a collocated $(r, z)$ grid with
$N_r$ radial and $N_z$ axial cells:
\begin{equation}\label{eq:grid}
  r_i = \bigl(i + \tfrac{1}{2}\bigr)\,\Delta r, \qquad
  z_j = \bigl(j + \tfrac{1}{2}\bigr)\,\Delta z
\end{equation}
where $\Delta r = R_{\mathrm{tube}} / N_r$ and
$\Delta z = L_{\mathrm{total}} / N_z$.  The total cell count is
$N_{\mathrm{cells}} = N_r \times N_z$.

\subsection{Flow Field Storage}

Each flow cell stores velocity $\bm{u}_{ij} = (u_r, u_\theta, u_z)$ and
pressure $p_{ij}$.  The flow field is initialised to the Hagen--Poiseuille
baseline:
\begin{equation}\label{eq:poiseuille_init}
  u_z(r_i) = 2\,\bar{U}\!\left(1 - \frac{r_i^2}{R_{\mathrm{tube}}^2}\right),
  \qquad
  p(z_j) = \Delta P_{\mathrm{HP}}\!\left(1 - \frac{j}{N_z}\right)
\end{equation}
where $\Delta P_{\mathrm{HP}} = 8\,\mu_f\,\bar{U}\,L_{\mathrm{total}} / R_{\mathrm{tube}}^2$.

The maximum velocity magnitude over the grid is
\begin{equation}\label{eq:umax}
  u_{\max} = \max_{i,j} |\bm{u}_{ij}|
\end{equation}

\subsection{Electric Field Storage and Finite-Difference Gradient}

Each electric cell stores potential $\phi_{ij}$, field
$\bm{E}_{ij}$, and local charge density $\rho_{e,ij}$.

The field is initialised to the coaxial analytical solution
(\S\ref{sec:phase3}, equations~\eqref{eq:helix_centre}--\eqref{eq:wire_volume}
are in Phase~1).  After each outer iteration the gradient is recomputed via
second-order central differences:
\begin{equation}\label{eq:E_fd}
  E_{r,ij} = -\frac{\phi_{i+1,j} - \phi_{i-1,j}}{2\,\Delta r},
  \qquad
  E_{z,ij} = -\frac{\phi_{i,j+1} - \phi_{i,j-1}}{2\,\Delta z}
\end{equation}
Interior cells only ($1 \le i \le N_r-2$, $1 \le j \le N_z-2$).

Peak and volume-averaged field magnitudes:
\begin{align}
  E_{\max} &= \max_{i,j} |\bm{E}_{ij}| \label{eq:Emax} \\
  \bar{E}  &= \frac{1}{|\Omega|}\sum_{i,j}|\bm{E}_{ij}|\,\Delta V_{ij}
             \approx \frac{1}{N_{\mathrm{cells}}}\sum_{i,j}|\bm{E}_{ij}|
             \label{eq:Eavg}
\end{align}

\subsection{Species Field and Flux Discretisation}

For each species~$i$ the concentration gradient at cell $(i_r, i_z)$ is
\begin{equation}\label{eq:grad_c}
  \nabla c_i\big|_{i_r, i_z} \approx
  \left(
    \frac{c_{i_r+1,i_z} - c_{i_r-1,i_z}}{2\,\Delta r},\;
    0,\;
    \frac{c_{i_r,i_z+1} - c_{i_r,i_z-1}}{2\,\Delta z}
  \right)
\end{equation}

The discrete Nernst--Planck flux (equation~\eqref{eq:NP}) at a cell is then
\begin{equation}\label{eq:NP_discrete}
  \bm{N}_{i}^{\,(k)} = -D_i\,\nabla c_i
                       - z_i\,\mu_i\,F\,c_i\,\nabla\phi
                       + c_i\,\bm{u}
\end{equation}
where all fields are evaluated at the $(k)$-th outer iteration.

\subsection{Charge-Density Assembly}

After species fluxes are updated, the charge density at each cell is
\begin{equation}\label{eq:rho_e_assembly}
  \rho_{e,ij} = F \sum_{s=1}^{N_s} z_s\,c_{s,ij}
\end{equation}
This feeds back into the Poisson solve for the next outer iteration.

\subsection{Segregated Outer-Iteration Algorithm}

The coupling loop (Algorithm~\ref{alg:coupling}) is:

\begin{enumerate}
  \item \textbf{Flow sub-solve}: update $\bm{u}, p$ from Navier--Stokes with
        current $\rho_e \bm{E}$ body force.  (Production: SIMPLE or projection
        method; baseline: retain Poiseuille.)
  \item \textbf{Electrostatic sub-solve}: recompute $\bm{E}$ from $\phi$ via
        \eqref{eq:E_fd}.  (Production: assemble $\rho_e$ into Poisson,
        solve with CG/GMRES.)
  \item \textbf{Transport sub-solve}: evaluate $\bm{N}_i$ for all species
        using current $\bm{u}$ and $\bm{E}$.  (Production: upwind advection
        + implicit diffusion.)
  \item \textbf{Charge update}: assemble $\rho_e$ from \eqref{eq:rho_e_assembly}.
  \item \textbf{Convergence check}: compute residual
        \begin{equation}\label{eq:residual}
          \mathcal{R} = \sqrt{\frac{1}{N_{\mathrm{cells}}}
                        \sum_{i,j}\rho_{e,ij}^2}
        \end{equation}
        Converged if $\mathcal{R} < \varepsilon_{\mathrm{tol}}$.
\end{enumerate}
\label{alg:coupling}

Under-relaxation is applied to each sub-solve:
\begin{equation}\label{eq:underrelax}
  \xi^{(k+1)} = \alpha\,\xi_{\mathrm{new}} + (1 - \alpha)\,\xi^{(k)}
\end{equation}
with factors $\alpha_{\mathrm{flow}}$, $\alpha_{\mathrm{elec}}$,
$\alpha_{\mathrm{species}}$ specified in the run card.

\subsection{Outlet Molar Flux Integration}

The species molar flow rate across the outlet is computed as
\begin{equation}\label{eq:outlet_flux}
  \dot{N}_i = \int_{\Gamma_{\mathrm{out}}} \bm{N}_i\cdot\bm{n}\,dA
  \approx \sum_{i_r=0}^{N_r-1} N_{i,z}\big|_{i_r,N_z-1}\;
           \cdot\; 2\pi\,r_{i_r}\,\Delta r
\end{equation}

\subsection{Accumulation Index}

A measure of species focusing or depletion by the electric field:
\begin{equation}\label{eq:accumulation}
  \mathcal{A}_i = \frac{\max_{i_r,i_z}\,c_{i}}
                       {\min_{i_r,i_z}\,c_{i}}
\end{equation}
$\mathcal{A}_i = 1$ for a perfectly uniform field; large values indicate
strong electromigration-driven concentration.

% ============================================================================
\section{Phase 4 — Mesh Sizing and Adaptive Refinement}
\label{sec:phase4}
% ============================================================================

Although the current solver uses a uniform structured grid, the mesh-control
framework specifies element-size fields that external meshers (gmsh, ANSYS
Meshing, snappyHexMesh) consume for unstructured or hybrid mesh generation.

\subsection{Global and Local Sizing}

A base element size $h_0$ is set proportional to the tube radius:
\begin{equation}\label{eq:h0}
  h_0 = \alpha_h\,R_{\mathrm{tube}}
\end{equation}
with $\alpha_h \approx 0.12$.  Local sizes are imposed near critical features:
\begin{align}
  h_{\mathrm{elec}} &= \beta_e\,d_w  & (\beta_e \approx 0.06) \label{eq:h_elec} \\
  h_{\mathrm{wall}} &= \beta_w\,R_{\mathrm{tube}} & (\beta_w \approx 0.025) \label{eq:h_wall}
\end{align}

A geometric growth rate~$\gamma$ transitions from the fine local size to the
global base size:
\begin{equation}\label{eq:growth}
  h_{n+1} = \gamma\,h_n, \qquad \gamma \approx 1.2
\end{equation}

\subsection{Boundary-Layer Inflation}

Near viscous walls and electrode surfaces, prism (inflation) layers resolve the
velocity and concentration boundary layers.  The first cell height is
\begin{equation}\label{eq:first_cell}
  \delta_1 = \alpha_\delta\,R_{\mathrm{tube}}
\end{equation}
with $\alpha_\delta \approx 0.003$.  Subsequent layers grow geometrically:
\begin{equation}\label{eq:inflation_growth}
  \delta_k = \delta_1\,\gamma_{\mathrm{inf}}^{\,k-1}, \qquad
  k = 1, \ldots, N_{\mathrm{layers}}
\end{equation}
Total inflation thickness must not exceed $\delta_{\max}$.

\subsection{Refinement Zones — Electrode Vicinity}

For each of the $n_t$ helix turns, a cylindrical refinement zone is centred
on the helix wire:
\begin{equation}\label{eq:refine_zone}
  \mathcal{Z}_t = \bigl\{(r, z) :
    |r - R_h| \le d_w,\;
    |z - z_{c,t}| \le 0.6\,p
  \bigr\}
\end{equation}
where $z_{c,t} = z_0 + (t + \tfrac{1}{2})\,p$ is the axial centre of turn~$t$.

The total number of refinement zones scales linearly with turn count:
\begin{equation}\label{eq:n_zones}
  N_{\mathrm{zones}} = n_t + N_{\mathrm{transition}}
\end{equation}
where $N_{\mathrm{transition}}$ accounts for inlet/outlet transition zones.

% ============================================================================
\section{Phase 5 — Postprocessing, Probes, and Comparison Framework}
\label{sec:phase5}
% ============================================================================

\subsection{Contour Field Export}

Scalar field quantities $\psi(r_i, z_j)$ are exported as structured $(r, z,
\psi)$ CSV grids.  The supported quantities and their definitions are:
\begin{align}
  |\bm{u}|_{ij}  &= \sqrt{u_{r,ij}^2 + u_{\theta,ij}^2 + u_{z,ij}^2}
                    & \text{(velocity magnitude)} \\
  p_{ij}          && \text{(pressure)} \\
  \phi_{ij}       && \text{(electric potential)} \\
  |\bm{E}|_{ij}  &= \sqrt{E_{r,ij}^2 + E_{z,ij}^2}
                    & \text{(field magnitude)} \\
  \rho_{e,ij}     && \text{(charge density)} \\
  c_{s,ij}        && \text{(species concentration)}
\end{align}

\subsection{Streamline Tracing}

Streamlines are integrated through the velocity field by forward-Euler
advection from inlet seed positions:
\begin{equation}\label{eq:euler_step}
  \bm{x}^{(n+1)} = \bm{x}^{(n)} + \Delta s\;\frac{\bm{u}(\bm{x}^{(n)})}{|\bm{u}(\bm{x}^{(n)})|}
\end{equation}
with step size $\Delta s = f_s\,\min(\Delta r, \Delta z)$ and step fraction
$f_s \approx 0.3$.  Seeds are distributed uniformly across the inlet radius:
\begin{equation}\label{eq:seeds}
  r_{0,k} = \frac{R_{\mathrm{tube}}(k + \tfrac{1}{2})}{N_{\mathrm{seeds}}},
  \qquad k = 0,\ldots,N_{\mathrm{seeds}}-1
\end{equation}

\subsection{Section Probes — Radial and Axial}

A \emph{radial probe} extracts the profile $\psi(r)$ at a fixed axial
station~$j^*$:
\begin{equation}\label{eq:radial_probe}
  \{(r_i, \psi_{ij^*})\}_{i=0}^{N_r-1}
\end{equation}

An \emph{axial probe} extracts $\psi(z)$ at a fixed radial index~$i^*$:
\begin{equation}\label{eq:axial_probe}
  \{(z_j, \psi_{i^*j})\}_{j=0}^{N_z-1}
\end{equation}

Both probes support all six field quantities listed in \S\ref{sec:phase5}.

\subsection{Comparison Tables and Enhancement Factor}

Given two solver metric sets $\mathcal{M}_{\mathrm{base}}$ (straight) and
$\mathcal{M}_{\mathrm{mod}}$ (modulated), the enhancement factor for each
metric is
\begin{equation}\label{eq:eta}
  \eta_m = \frac{m_{\mathrm{mod}}}{m_{\mathrm{base}}},
  \qquad m \in \{\Delta P,\; E_{\max},\; \bar{E},\;
  u_{\max},\; \mathrm{Re},\; \dot{N}_{i},\; \mathcal{A}_i\}
\end{equation}

\subsection{Multi-Objective Optimisation}

The weighted objective function for geometry optimisation is
\begin{equation}\label{eq:objective}
  J(\mathcal{P}) = w_1\,\bar{E} + w_2\,\dot{N}_{\mathrm{out}} - w_3\,\Delta P
\end{equation}
The negative sign on $\Delta P$ penalises excessive pressure loss.  Pareto
exploration sweeps weight vectors $\bm{w} = (w_1, w_2, w_3)$ to identify
non-dominated geometry candidates.

% ============================================================================
\section{Phase 6 — Dimensional Analysis, Sensitivity, and Verification}
\label{sec:phase6}
% ============================================================================

\subsection{Dimensionless Groups}

The EHD system is governed by the following key dimensionless numbers:

\begin{align}
  \mathrm{Re}   &= \frac{\rho\,\bar{U}\,D_h}{\mu_f}
                  & \text{(Reynolds — inertia/viscous)}
                  \label{eq:Re_dim} \\[4pt]
  \mathrm{Eh}   &= \frac{\epsilon\,\Delta V}{\mu_f\,\bar{U}\,D_h}
                  & \text{(electric Hartmann analogue)}
                  \label{eq:Eh} \\[4pt]
  \mathrm{Pe}_i &= \frac{\bar{U}\,D_h}{D_i}
                  & \text{(Péclet — convection/diffusion)}
                  \label{eq:Pe} \\[4pt]
  \mathrm{Da}_i &= \frac{|z_i|\,\mu_i\,F\,\Delta V}{D_i}
                  & \text{(Damköhler-migration — migration/diffusion)}
                  \label{eq:Da}
\end{align}

These groups classify the dominant physics regime: for $\mathrm{Re} \ll 2300$
the flow is laminar; for $\mathrm{Pe}_i \gg 1$ convection dominates species
transport; for $\mathrm{Da}_i \gg 1$ electromigration dominates diffusion.

\subsection{Parametric Sensitivity — Voltage Sweep}

At constant flow rate, increasing $\Delta V$ scales the electric field
linearly:
\begin{equation}\label{eq:E_vs_V}
  E_{\max}(\Delta V) = \frac{\Delta V}{R_{\mathrm{in}}\,
    \ln(R_{\mathrm{out}} / R_{\mathrm{in}})}
  \propto \Delta V
\end{equation}
while the pressure drop remains approximately constant (the Lorentz body
force is small compared to the viscous force for typical EHD operating
conditions at low $\mathrm{Re}$).

\subsection{Parametric Sensitivity — Velocity Sweep}

The analytical pressure drop scales linearly with mean velocity:
\begin{equation}\label{eq:dP_vs_U}
  \Delta P = \frac{8\,\mu_f\,\bar{U}\,L_{\mathrm{total}}}
                  {R_{\mathrm{tube}}^2}
  \propto \bar{U}
\end{equation}
and the Reynolds number also scales linearly:
\begin{equation}\label{eq:Re_vs_U}
  \mathrm{Re} = \frac{\rho\,\bar{U}\,D_h}{\mu_f} \propto \bar{U}
\end{equation}

\subsection{Memory Footprint Estimation}

Per-cell storage requirements:
\begin{align}
  S_{\mathrm{flow}} &= |\texttt{FlowCell}| = 32\;\text{bytes}
    & (3\times 8\;\text{B velocity} + 8\;\text{B pressure}) \\
  S_{\mathrm{elec}} &= |\texttt{ElectricCell}| = 40\;\text{bytes}
    & (8\;\text{B }\phi + 3\times 8\;\text{B }\bm{E} + 8\;\text{B }\rho_e) \\
  S_{\mathrm{spec}} &= |\texttt{SpeciesCell}| = 32\;\text{bytes}
    & (8\;\text{B }c_i + 3\times 8\;\text{B }\bm{N}_i)
\end{align}
Total memory for $N_s$ species on an $N_r \times N_z$ grid:
\begin{equation}\label{eq:memory}
  M_{\mathrm{total}} = N_r\,N_z\,\bigl(S_{\mathrm{flow}}
                      + S_{\mathrm{elec}}
                      + N_s\,S_{\mathrm{spec}}\bigr)
\end{equation}

For the default single-species case ($N_s = 1$) on a $160\times 800$ grid
($128{,}000$ cells):
\begin{equation}
  M \approx 128{,}000 \times (32 + 40 + 32) = 13.31\;\text{MB}
\end{equation}

\subsection{Determinism Verification}

A core design requirement of the VSEPR-SIM platform is bit-exact
reproducibility.  All field computations are IEEE~754 double-precision with
deterministic reduction order.  Verification confirms that repeated
execution produces identical output metrics:
\begin{equation}\label{eq:determinism}
  \bigl|\,\Delta P^{(a)} - \Delta P^{(b)}\bigr| = 0,
  \qquad
  \bigl|\,E_{\max}^{(a)} - E_{\max}^{(b)}\bigr| = 0,
  \qquad
  \bigl|\,u_{\max}^{(a)} - u_{\max}^{(b)}\bigr| = 0
\end{equation}
for any two invocations $(a)$, $(b)$ with identical input parameters.

% ============================================================================
\section{Phase 7 — EHD Pump Configurations and Pumping Mechanisms}
\label{sec:phase7}
% ============================================================================

The preceding phases treated a single device topology (helical wire electrode
inside a cylindrical tube).  Real EHD pump systems span a broader design space
with at least four principal electrode configurations and three distinct
pumping mechanisms.  This phase formalises the theoretical basis for each.

% ----------------------------------------------------------------------------
\subsection{Four Principal Configurations}
% ----------------------------------------------------------------------------

\subsubsection{Configuration (a): Planar Electrode Channel}

Two parallel planar electrodes separated by a gap~$H$ drive flow through a
rectangular fluidic channel of width~$W$ and streamwise length~$L_{ch}$.

\paragraph{Uniform field.}
The electric field between infinite parallel plates is
\begin{equation}\label{eq:planar_field}
  E_y = \frac{\Delta V}{H}
\end{equation}
and the potential varies linearly:
\begin{equation}\label{eq:planar_phi}
  \phi(y) = \Delta V\!\left(1 - \frac{y}{H}\right)
\end{equation}

\paragraph{Plane-Poiseuille baseline.}
The velocity profile between no-slip plates is
\begin{equation}\label{eq:plane_poiseuille}
  u_z(y) = \frac{6\,\bar{U}}{H^2}\,y\,(H - y)
\end{equation}
with peak velocity $u_{\max} = \tfrac{3}{2}\bar{U}$ at the midplane.  The
associated pressure drop is
\begin{equation}\label{eq:plane_dp}
  \Delta P = \frac{12\,\mu_f\,\bar{U}\,L_{ch}}{H^2}
\end{equation}

\paragraph{Hydraulic diameter.}
\begin{equation}\label{eq:D_h_rect}
  D_h = \frac{2\,H\,W}{H + W}
\end{equation}

\subsubsection{Configuration (b): Needle-Ring Ion Injection}

A sharp needle electrode (tip curvature~$r_{\mathrm{tip}}$) faces a grounded
ring electrode (inner radius~$R_{\mathrm{ring}}$) across an axial
gap~$d_{\mathrm{gap}}$.  Corona discharge at the needle tip injects ions that
are accelerated by the field, transferring momentum to the neutral gas/liquid.

\paragraph{Tip field enhancement.}
For a hyperboloidal needle the field at the tip is
\begin{equation}\label{eq:needle_tip}
  E_{\mathrm{tip}} \approx
    \frac{\Delta V}{r_{\mathrm{tip}}
    \,\ln\!\bigl(2\,d_{\mathrm{gap}} / r_{\mathrm{tip}}\bigr)}
\end{equation}
The ratio $E_{\mathrm{tip}} / ({\Delta V}/{d_{\mathrm{gap}}})$ is the
geometric enhancement factor, typically $\gg 1$ for sharp tips.

\paragraph{Peek's corona onset criterion.}
In air at relative density~$\delta$ (STP $\delta = 1$) the onset field for
positive corona is
\begin{equation}\label{eq:peek}
  E_{\mathrm{onset}} = 3.1\times10^{6}\;\delta\!
    \left(1 + \frac{0.308}{\sqrt{\delta\,r_{\mathrm{tip}}}}\right)
  \quad [\text{V/m}]
\end{equation}

\paragraph{Ion-drag pressure.}
The generated pressure head from corona current~$I$ is
\begin{equation}\label{eq:ion_drag_P}
  \Delta P_{\mathrm{ion}} \approx
    \frac{I\,d_{\mathrm{gap}}}{\mu_i\,A_{\mathrm{ring}}}
\end{equation}
where $\mu_i$ is the ion mobility and
$A_{\mathrm{ring}} = \pi R_{\mathrm{ring}}^2$.

\subsubsection{Configuration (c): Stacked Perforated Disk Electrodes}

$N_d$ circular disks of radius~$R_{\mathrm{disk}}$ and thickness~$t_d$
are arranged along the flow axis with alternating polarity.  Each disk is
perforated with $N_p$ through-holes of radius~$r_p$.

\paragraph{Stack geometry.}
Total stack length:
\begin{equation}\label{eq:stack_length}
  L_{\mathrm{stack}} = N_d\,t_d + (N_d - 1)\,d_s
\end{equation}
where $d_s$ is the inter-disk spacing.  The open-area fraction is
\begin{equation}\label{eq:oaf}
  \alpha_{\mathrm{OA}} = \frac{N_p\,\pi\,r_p^2}{\pi\,R_{\mathrm{disk}}^2}
  = \frac{N_p\,r_p^2}{R_{\mathrm{disk}}^2}
\end{equation}

\paragraph{Inter-disk field.}
Between adjacent disks (uniform approximation):
\begin{equation}\label{eq:disk_field}
  E_{\mathrm{gap}} = \frac{\Delta V}{d_s}
\end{equation}

\paragraph{Flow character.}
The device produces both axial through-flow (via perforations) and rotational
swirl in the inter-disk gaps.  The effective cross-section for through-flow is
$A_{\mathrm{eff}} = N_p\,\pi\,r_p^2$.

\subsubsection{Configuration (d): Triangular Prism Slit-Electrode}

An array of $N_{\mathrm{prism}}$ triangular prism electrodes with base~$b$
and height~$h_p$ are separated by slit gaps of width~$w_s$.  The prism tips
produce strongly non-uniform fields that drive cross-flow patterns.

\paragraph{Array pitch.}
\begin{equation}\label{eq:prism_pitch}
  p_{\mathrm{array}} = b + w_s
\end{equation}
Total array width:
$W_{\mathrm{array}} = N_{\mathrm{prism}}\,b + (N_{\mathrm{prism}} - 1)\,w_s$.

\paragraph{Tip field enhancement.}
For an isosceles triangle with half-angle
$\alpha = \arctan(b/2h_p)$, the field at the apex is enhanced relative to
the uniform slit field:
\begin{equation}\label{eq:prism_tip}
  E_{\mathrm{tip}} \approx \frac{\Delta V}{w_s} \cdot \frac{1}{\sin\alpha}
\end{equation}
The factor $\beta = 1/\sin\alpha$ captures the lightning-rod concentration
of field lines at sharp tips.

% ----------------------------------------------------------------------------
\subsection{Three EHD Pumping Mechanisms}
% ----------------------------------------------------------------------------

\subsubsection{Mechanism 1: Coulomb / Ion-Drag Force}

The dominant mechanism in ion-injection and corona-discharge devices.
Space charge $\rho_e$ interacts with the local field:
\begin{equation}\label{eq:f_coulomb}
  \bm{f}_{\mathrm{C}} = \rho_e\,\bm{E}
\end{equation}
This produces a volumetric body force that is injected into the momentum
equation.  The characteristic Coulomb pressure scale is
\begin{equation}\label{eq:P_coulomb}
  P_C \sim \epsilon\,E^2 = \epsilon\!\left(\frac{\Delta V}{d}\right)^{\!2}
\end{equation}

\subsubsection{Mechanism 2: Dielectrophoretic (DEP) Force}

Arises from permittivity gradients in non-uniform electric fields.  The
time-averaged DEP force on a polarisable medium is
\begin{equation}\label{eq:f_dep}
  \bm{f}_{\mathrm{DEP}} = \frac{1}{2}\,\epsilon_0\,K_{\mathrm{CM}}
    \,\nabla|\bm{E}|^2
\end{equation}
where $K_{\mathrm{CM}} = \mathrm{Re}[K(\omega)]$ is the Clausius-Mossotti
factor (bounded between $-0.5$ and $+1.0$ for spherical particles).

The characteristic DEP velocity is
\begin{equation}\label{eq:u_dep}
  u_{\mathrm{DEP}} \sim
    \frac{\epsilon\,K_{\mathrm{CM}}\,E^2\,d}{\mu_f}
\end{equation}

\subsubsection{Mechanism 3: Electroosmotic Flow (EOF)}

In micro-channels with charged walls, the electric double layer (EDL)
interacts with the tangential field to produce a plug-flow slip at the wall.

\paragraph{Helmholtz-Smoluchowski slip.}
\begin{equation}\label{eq:eof_slip}
  \bm{u}_{\mathrm{slip}} = -\frac{\epsilon\,\zeta}{\mu_f}\,\bm{E}_t
\end{equation}
where $\zeta$ is the wall zeta potential and $\bm{E}_t$ is the
tangential (wall-parallel) component of the field.

\paragraph{Debye length.}
The characteristic thickness of the EDL is
\begin{equation}\label{eq:debye}
  \lambda_D = \sqrt{\frac{\epsilon\,R\,T}{2\,F^2\,c_0}}
\end{equation}
The thin-EDL approximation ($\kappa\,H \gg 1$ where $\kappa = 1/\lambda_D$)
is valid when the Debye length is much smaller than the channel dimension.

\paragraph{EOF flow rate (capillary).}
In a cylindrical capillary of radius~$R_c$:
\begin{equation}\label{eq:Q_eof}
  Q_{\mathrm{EOF}} = -\frac{\epsilon\,\zeta\,E_z\,\pi\,R_c^2}{\mu_f}
\end{equation}

\paragraph{EOF pressure head.}
Under closed-channel conditions the electroosmotic flow generates a
back-pressure:
\begin{align}
  \Delta P_{\mathrm{EOF,cyl}}    &= \frac{8\,\epsilon\,|\zeta|\,|E_z|\,L}{R_c^2}
    \label{eq:dP_eof_cyl} \\
  \Delta P_{\mathrm{EOF,planar}} &= \frac{12\,\epsilon\,|\zeta|\,|E_z|\,L}{H^2}
    \label{eq:dP_eof_planar}
\end{align}

% ----------------------------------------------------------------------------
\subsection{Topology--Mechanism Affinity}
% ----------------------------------------------------------------------------

Not every mechanism is equally relevant to every configuration.  The natural
pairings are:

\begin{table}[h]
\centering
\caption{Primary pumping mechanisms by device topology}
\begin{tabular}{@{}lcccp{4.5cm}@{}}
\toprule
Configuration & Coulomb & DEP & EOF & Typical Application \\
\midrule
(a) Planar channel    & \checkmark & --- & \checkmark & Micro-pump, lab-on-chip \\
(b) Needle-ring       & \checkmark & --- & ---        & Gas pump, ion wind \\
(c) Disk stack        & \checkmark & --- & ---        & Multi-stage liquid pump \\
(d) Prism slit        & \checkmark & \checkmark & --- & Particle manipulation \\
Helical wire (orig.)  & \checkmark & --- & ---        & Tubular reactor \\
\bottomrule
\end{tabular}
\end{table}

% ############################################################################
%   PHASE 8 — Metallic Particle Combustion and Reactive Multiphase Flow
% ############################################################################
\section{Phase 8: Metallic Particle Combustion and Reactive Multiphase Flow}
\label{sec:phase8}

This phase introduces a Lagrangian--Eulerian reactive multiphase model
coupling discrete burning particles with the carrier-gas continuum.  Five
representative fuel types—methane~(CH$_4$), iron~(Fe), aluminum~(Al),
boron/aluminum~(B/Al), and zirconium~(Zr)—are modelled following the
experimental observations of Bergthorson (2015) and the classical combustion
theory of Glassman \& Yetter.

% ----------------------------------------------------------------------------
\subsection{Fuel Classification}
% ----------------------------------------------------------------------------

The five fuels span three combustion regimes:

\begin{table}[h]
\centering
\caption{Fuel types, combustion regimes, and characteristic emission species}
\begin{tabular}{@{}llllr@{}}
\toprule
Fuel & Formula & Regime & Emitter & $T_{\mathrm{ad}}$ (K) \\
\midrule
Methane        & CH$_4$  & Premixed gas  & CH$^*$  & 2\,223 \\
Iron           & Fe      & Surface       & FeO     & 2\,340 \\
Aluminum       & Al      & Vapour-phase  & AlO     & 3\,673 \\
Boron/Aluminum & B/Al    & Hybrid        & BO$_2$  & 3\,450 \\
Zirconium      & Zr      & Hybrid        & ZrO     & 2\,950 \\
\bottomrule
\end{tabular}
\end{table}

% ----------------------------------------------------------------------------
\subsection{d\textsuperscript{n}-Law Burning Model}
% ----------------------------------------------------------------------------

The classical shrinking-particle model generalises the d$^2$-law to an
arbitrary exponent~$n$:
\begin{equation}\label{eq:dn_law}
  d(t)^{n} = d_0^{n} - K\,t
\end{equation}
where $K$~(m$^n$/s) is the \emph{burning-rate constant}.  For
vapour-phase combustion (Al, B/Al) the classical $n=2$ applies; for
surface-limited heterogeneous oxidation (Fe) $n\approx 1$; and for hybrid
fuels (Zr) $n\approx 1.8$.

The particle burnout time follows immediately:
\begin{equation}\label{eq:t_burn}
  t_{\mathrm{burn}} = \frac{d_0^{n}}{K}
\end{equation}

The instantaneous mass-loss rate for a spherical particle of density
$\rho_p$ is:
\begin{equation}\label{eq:dm_dt}
  \frac{\mathrm{d}m}{\mathrm{d}t}
    = -\frac{\pi}{2}\,\rho_p\,K\,d^{\,2-n}
\end{equation}
which is \emph{constant} for $n=2$ (classical d$^2$-law) and linearly
proportional to~$d$ for $n=1$ (surface regime).

% ----------------------------------------------------------------------------
\subsection{Burning-Rate Corrections}
% ----------------------------------------------------------------------------

\paragraph{Temperature correction (Arrhenius).}
\begin{equation}\label{eq:K_arrhenius}
  K(T) = K_0 \exp\!\left[\frac{E_a}{R}
         \left(\frac{1}{T_0} - \frac{1}{T}\right)\right]
\end{equation}
where $E_a$ is the activation energy and $T_0=300$\,K is the reference
temperature.

\paragraph{Oxygen fraction correction.}
\begin{equation}\label{eq:K_O2}
  K(Y_{\mathrm{O}_2}) = K_0 \,\frac{Y_{\mathrm{O}_2}}{0.233}
\end{equation}

\paragraph{EHD electric-field enhancement.}
An applied electric field modifies combustion via ionic-wind-enhanced
transport (Lawton \& Weinberg):
\begin{equation}\label{eq:K_ehd}
  K_E = K_0 \left[1 + \alpha_E
        \left(\frac{E}{E_{\mathrm{ref}}}\right)^{\!2}\right]
\end{equation}
with $\alpha_E \approx 0.05$--$0.3$ (fuel-dependent) and
$E_{\mathrm{ref}} = 10$\,kV/m.

% ----------------------------------------------------------------------------
\subsection{Heat Release}
% ----------------------------------------------------------------------------

The instantaneous heat release rate from a single particle is:
\begin{equation}\label{eq:Q_dot}
  \dot{Q} = \left|\frac{\mathrm{d}m}{\mathrm{d}t}\right|\,\Delta H_c
\end{equation}

The adiabatic flame temperature at stoichiometric conditions in air:
\begin{equation}\label{eq:T_ad}
  T_{\mathrm{ad}} = T_0
    + \frac{\Delta H_c}{c_{p,\mathrm{mix}}\,(1 + \mathrm{AF}_{\mathrm{stoich}})}
\end{equation}

% ----------------------------------------------------------------------------
\subsection{Radiative Emission}
% ----------------------------------------------------------------------------

The Stefan--Boltzmann radiative power from a spherical particle of
diameter~$d$ at temperature~$T_p$:
\begin{equation}\label{eq:P_rad}
  P_{\mathrm{rad}} = \varepsilon_p\,\sigma\,\pi\,d^2\,
    \left(T_p^4 - T_\infty^4\right)
\end{equation}

The distinct flame colours observed in Figure~Z.3 arise from
element-specific electronic transitions.  The peak emission wavelength
follows Wien's displacement law:
\begin{equation}\label{eq:wien}
  \lambda_{\mathrm{peak}} = \frac{b}{T_p}
  \qquad (b = 2.898 \times 10^{-3}~\mathrm{m \cdot K})
\end{equation}

The spectral radiance at wavelength~$\lambda$ is given by Planck's law:
\begin{equation}\label{eq:planck}
  B(\lambda, T) = \frac{2hc^2}{\lambda^5}
    \,\frac{1}{\exp\!\left(\dfrac{hc}{\lambda\,k_B\,T}\right) - 1}
\end{equation}

% ----------------------------------------------------------------------------
\subsection{Lagrangian Particle Phase}
% ----------------------------------------------------------------------------

Each particle is tracked as a Lagrangian entity with state
$(x_p, v_p, d, T_p, m)$.  The equations of motion are:

\paragraph{Position.}
\begin{equation}\label{eq:dxp_dt}
  \frac{\mathrm{d}x_p}{\mathrm{d}t} = v_p
\end{equation}

\paragraph{Velocity (Stokes drag + EHD force).}
\begin{equation}\label{eq:dvp_dt}
  \frac{\mathrm{d}v_p}{\mathrm{d}t}
    = \frac{u - v_p}{\tau_p} + \frac{F_E}{m_p}
\end{equation}
where the particle relaxation time is:
\begin{equation}\label{eq:tau_p}
  \tau_p = \frac{\rho_p\,d^2}{18\,\mu_g}
\end{equation}
and $F_E = q_p\,E$ is the Lorentz force on a thermionically charged
particle.

\paragraph{Temperature (heat balance).}
\begin{equation}\label{eq:dTp_dt}
  m_p\,c_p\,\frac{\mathrm{d}T_p}{\mathrm{d}t}
    = \dot{Q}_{\mathrm{gen}} - P_{\mathrm{rad}}
\end{equation}

\paragraph{Thermionic charge.}
Burning metal particles acquire charge via thermionic emission.
The Rayleigh-limit estimate gives:
\begin{equation}\label{eq:q_p}
  q_p \approx 4\pi\varepsilon_0\,r_p\,\frac{k_B\,T_p}{e}
\end{equation}

% ----------------------------------------------------------------------------
\subsection{Eulerian Source Terms (Two-Way Coupling)}
% ----------------------------------------------------------------------------

Particle effects are projected onto the Eulerian grid as volumetric
source terms.  For $N_p$ particles in a cell of volume~$V_{\mathrm{cell}}$:

\paragraph{Mass source.}
\begin{equation}\label{eq:S_mass}
  S_{\mathrm{mass}} = \frac{N_p}{V_{\mathrm{cell}}}
    \left|\frac{\mathrm{d}m}{\mathrm{d}t}\right|
\end{equation}

\paragraph{Energy source.}
\begin{equation}\label{eq:S_energy}
  S_{\mathrm{energy}} = \frac{N_p}{V_{\mathrm{cell}}}\,\dot{Q}
\end{equation}

\paragraph{Momentum source (drag coupling).}
\begin{equation}\label{eq:S_mom}
  S_{\mathrm{mom}} = \frac{18\,\mu_g\,N_p}{\rho_p\,d^2\,V_{\mathrm{cell}}}
    \,(u - v_p)
\end{equation}

\paragraph{Oxygen consumption.}
\begin{equation}\label{eq:S_O2}
  S_{\mathrm{O}_2} = -\frac{N_p}{V_{\mathrm{cell}}}
    \left|\frac{\mathrm{d}m}{\mathrm{d}t}\right|
    \frac{0.233\,\mathrm{AF}_{\mathrm{stoich}}}{1 + \mathrm{AF}_{\mathrm{stoich}}}
\end{equation}

% ----------------------------------------------------------------------------
\subsection{Dimensionless Numbers}
% ----------------------------------------------------------------------------

\paragraph{Damk\"ohler number (first kind).}
The ratio of flow residence time to reaction burnout time:
\begin{equation}\label{eq:Da_I}
  \mathrm{Da}_I = \frac{\tau_{\mathrm{flow}}}{\tau_{\mathrm{rxn}}}
    = \frac{L / U}{d_0^{n} / K}
\end{equation}
$\mathrm{Da} \gg 1$: fast reaction (mixing-limited);
$\mathrm{Da} \ll 1$: slow reaction (kinetics-limited).

\paragraph{Stokes number.}
\begin{equation}\label{eq:St_comb}
  \mathrm{St} = \frac{\tau_p}{\tau_f}
    = \frac{\rho_p\,d^2}{18\,\mu_g} \cdot \frac{U}{L}
\end{equation}

\paragraph{Mass loading ratio.}
\begin{equation}\label{eq:phi_m}
  \phi_m = \frac{N_p\,m_p}{\rho_g\,Q}
\end{equation}
$\phi_m < 0.1$: one-way coupling; $\phi_m > 0.1$: two-way coupling required.

\paragraph{Biot number.}
\begin{equation}\label{eq:Bi}
  \mathrm{Bi} = \frac{h\,d}{2\,k_p}
\end{equation}
For metal particles $\mathrm{Bi} \ll 1$ (lumped-capacitance valid).

\paragraph{Nusselt number (Ranz--Marshall).}
\begin{equation}\label{eq:Nu_RM}
  \mathrm{Nu} = 2 + 0.6\,\mathrm{Re}_p^{1/2}\,\mathrm{Pr}^{1/3}
\end{equation}

\paragraph{Particle Reynolds number.}
\begin{equation}\label{eq:Re_p}
  \mathrm{Re}_p = \frac{\rho_g\,|u - v_p|\,d}{\mu_g}
\end{equation}

% ----------------------------------------------------------------------------
\subsection{Fuel Comparison Framework}
% ----------------------------------------------------------------------------

The five fuels are compared using:
\begin{itemize}
  \item Gravimetric energy density: $e_g = \Delta H_c$ (J/kg)
  \item Volumetric energy density: $e_v = \rho\,\Delta H_c$ (J/m$^3$)
  \item Burnout time for $d_0 = 10\,\mu$m reference particle
  \item Adiabatic flame temperature $T_{\mathrm{ad}}$
  \item Emissivity and peak emission wavelength
\end{itemize}

\begin{table}[h]
\centering
\caption{Fuel comparison for $d_0 = 10\,\mu$m particles}
\begin{tabular}{@{}lrrrrl@{}}
\toprule
Fuel & $e_g$ (MJ/kg) & $e_v$ (GJ/m$^3$) & $K$ (m$^n$/s) & $t_b$ ($\mu$s) & $\varepsilon$ \\
\midrule
CH$_4$ & 50.0  & 0.033  & ---         & ---   & 0.10 \\
Fe     &  7.4  & 58.3   & $3.0\!\times\!10^{-7}$ & 33   & 0.85 \\
Al     & 31.1  & 84.0   & $2.0\!\times\!10^{-6}$ & 50   & 0.95 \\
B/Al   & 40.0  & 100.0  & $1.5\!\times\!10^{-6}$ & 67   & 0.88 \\
Zr     & 12.0  & 78.1   & $4.0\!\times\!10^{-7}$ & 68   & 0.80 \\
\bottomrule
\end{tabular}
\end{table}

% ============================================================================
\section*{Equation Index}
% ============================================================================

This document contains numbered equations grouped across eight expanded phases
and the preceding foundational sections.  Equation numbering is sequential and
cross-referenceable via \texttt{\textbackslash eqref}.

\end{document}
